\documentclass[12pt]{article}
\usepackage[utf8]{inputenc}
\usepackage[T1]{fontenc}
\usepackage[margin=1in]{geometry}
\usepackage{setspace}
\usepackage{xcolor}
\usepackage[colorlinks=true, urlcolor=blue!60!black, linkcolor=blue!60!black]{hyperref}
\setlength{\parindent}{2em}
\setlength{\parskip}{1.25em}
\renewcommand{\baselinestretch}{1.0}
\begin{document}
\begin{center}
{\LARGE \textit{Wolford v. Lopez}, 609 U.S. 185 (2026)\par}
\vspace{0.8em}
{\large Case Brief --- Four-Page Edition\par}
\end{center}
\vspace{1em}

\section{Memory Jogger}

In a 6--3 decision, the Court struck down Hawaii's post-\textit{Bruen} ``vampire rule'' barring licensed gun owners from carrying handguns onto private property open to the public without the owner's express permission, holding that states cannot flip the common-law implied-license default to hobble public carry (\href{https://www.supremecourt.gov/opinions/25pdf/24-1046_nmio.pdf}{Slip Op.}).

\section{Detailed Case Facts}

Before 2022 Hawaii issued carry permits only in ``exceptional case[s],'' reportedly just four between 2000 and 2018. After \textit{New York State Rifle \& Pistol Ass'n v. Bruen} (2022) recognized a right to carry handguns outside the home and singled out Hawaii, the legislature enacted Act 52 (2023). It imposed rigorous licensing, banned guns in many ``sensitive places,'' and, in \S{}134-9.5, barred licensed carriers from bringing any firearm onto ``private property of another'' without ``express authorization'' by signage or ``[u]nambiguous written or verbal authorization.''

Critics dubbed \S{}134-9.5 the ``vampire rule,'' after the folklore that a vampire cannot cross a threshold uninvited. It reversed the common-law default under which opening property to the public creates an implied license to enter unless the owner revokes consent. California, Maryland, New Jersey, and New York adopted similar rules. Three Maui County permit holders, Jason and Alison Wolford and Atom Kasprzycki, with the Hawaii Firearms Coalition, sued Attorney General Anne E. Lopez.

\section{Procedural History}

The District of Hawaii preliminarily enjoined \S{}134-9.5 as applied to property open to the public. A Ninth Circuit panel reversed, finding a historical tradition allowing states to set consent defaults for armed entry, and the full court denied rehearing en banc over dissents. The Supreme Court granted review limited to property open to the public, heard argument on January 20, 2026, and reversed on June 25, 2026.

\section{Judicial Votes}

Reversed, \textbf{6--3}.
\begin{itemize}
\item \textbf{Majority:} Justice Alito, joined by Chief Justice Roberts and Justices Thomas, Gorsuch, Kavanaugh, and Barrett.
\item \textbf{Concurrence:} Justice Barrett, joined by Justices Thomas and Gorsuch as to Part II--B.
\item \textbf{Dissents:} Justice Kagan; Justice Jackson, joined by Justice Sotomayor.
\end{itemize}

\section{Holding}

A state may not bar licensed carriers from bringing handguns onto private property open to the public absent the owner's express authorization. Such an opt-in default falls within the Second Amendment's plain text at \textit{Bruen} step one and lacks a ``relevantly similar'' historical analogue at step two. Property-law limits belong at step two; local customs cannot dilute the right's national scope; and colonial anti-poaching laws, an 1893 Oregon statute, and an 1865 Louisiana Black Code are not adequate analogues. The Court did not disturb owners' right to exclude, property closed to the public, or the sensitive-place bans, and left open whether 1791 or 1868 controls.

\section{Analysis of Opinions}

\textbf{Alito (majority).} Petitioners are ``the people'' seeking to ``bear'' ``Arms.'' A carrier modeled on Jaime Caetano running errands to a gas station, restaurant, drug store, and grocery could become ``a criminal at least six times over,'' since few owners post ``Guns Welcome'' signs. Under \textit{McDonald}, the right has a uniform national meaning. The colonial laws targeted poaching on rural or enclosed lands, the Oregon law was an outlier, and the Louisiana statute was a ``tainted artifact'' of the Black Codes.

\textbf{Barrett (concurring).} The Second Amendment binds state action, so a state property law aimed at arms-bearing triggers scrutiny, as a permission rule for wearing a hijab would under the First Amendment. The founding-era laws support focused defaults for places prone to particular abuses, but Hawaii targeted none. Joined by Thomas and Gorsuch, she found the Black Codes' purpose of subordinating freedmen not relevantly similar.

\textbf{Kagan (dissenting).} At step two the ``how'' is identical and the ``why'' close enough under \textit{Rahimi}.

\textbf{Jackson (dissenting).} The case is ``about property rights, not gun rights.'' Because there is no right to carry onto another's property without consent, rules about the form of consent do not implicate the Second Amendment. She accused the majority of ``mov[ing] the goal posts'' and creating a ``one-way ratchet''; demanding a closer match to poaching traps the law ``in amber,'' and the Court must confront the record ``warts and all.''

\section{Examples: Future Applications}

\begin{enumerate}
\item \textbf{Same-Side (Invalid): Signage-Only Rule.} A state bars carry into any business without a standardized ``Firearms Permitted'' placard. It is stricter than Hawaii's rule and flips the same default.
\item \textbf{Same-Side (Invalid): Parking Lots.} A city bans guns, even locked in cars, on retail parking lots unless the owner files an annual opt-in registration. It flips the implied license without a historical analogue.
\item \textbf{Opposite-Side (Valid): ``No Firearms'' Signs.} A grocery chain posts signs and the state enforces general trespass law against a carrier who refuses to leave. Owners remain free to exclude.
\item \textbf{Opposite-Side (Valid): Homes and Farms.} A state requires express consent to carry onto residences and non-public farmland. Such property carries no implied license and parallels colonial enclosed-land laws.
\item \textbf{Fence-Sitter: Nightclubs and Dispensaries.} After late-night shootings, a state requires express consent only at nightclubs and cannabis dispensaries. Alito's opinion cuts against it, but Barrett's ``focused regulation'' principle may support it.
\end{enumerate}

\section{Critique}

\textbf{Originalist and textualist views.} Joseph Greenlee and George Mocsary praise the decision for blocking post-\textit{Bruen} evasion that would turn a right into a permission slip, while preserving owners' choices. Even sympathetic originalists worry that framing colonial laws as ``anti-poaching'' revives the ``historical twin'' rigidity \textit{Rahimi} disavowed.

\textbf{Progressive and property-law views.} Maureen Brady highlights the problem of choosing a level of generality, and the \textit{Harvard Law Review} warns that excluding the Black Code creates an unguided ``\textit{Bruen} step 1.5'' for judging tainted history. Everytown argues the decision overrides state federalism and burdens businesses wary of posting signs.

\section{Key Quotations}

\begin{quote}
``This regime hobbles what the Second Amendment protects: the right of Americans to carry arms for self-defense as they go about their daily lives. We hold that the law is unconstitutional.''
\end{quote}
\noindent --- Justice Alito, Opinion of the Court.

\begin{quote}
``But the Second Amendment does not apply to private parties. It does apply to the States. See U. S. Const., Amdt. 14, \S{}1. And when a State enacts a property law that regulates arms-bearing conduct, that law implicates the Second Amendment.''
\end{quote}
\noindent --- Justice Barrett, concurring.

\begin{quote}
``With this decision, the Court has now manipulated \textit{Bruen} into a free-for-all that lets the Judiciary thwart the will of legislatures by privileging access to firearms above all else. Today's decision makes one thing clear: The Court's objective is protecting guns, not consistently preserving any principle of law.''
\end{quote}
\noindent --- Justice Jackson, dissenting.

\end{document}
